\documentclass[../../main.tex]{subfiles}
\begin{document}
% SOURCE p=256; printed=241; object=ch08-conjugate-convolution; number=8.8; region=upper-middle
\begin{proposition}\label{proposition:ch08-conjugate-convolution}
若 $p,q$ 为共轭指数，$f\in L^p$、$g\in L^q$，则 $f*g(x)$ 对每个 $x$ 存在，$f*g$ 有界且一致连续，并有
\begin{align*}
\|f*g\|_u\leqslant\|f\|_p\|g\|_q.
\end{align*}
若 $1<p<\infty$（从而也有 $1<q<\infty$），则 $f*g\in C_0(\mathbb R^n)$。
\end{proposition}
\begin{proof}
存在性及一致范数估计直接来自 Hölder 不等式。由\cref{proposition:ch08-lp-translation-continuity,proposition:ch08-convolution-properties}，一致连续性也随之得到：若 $1\leqslant p<\infty$，则
\begin{align*}
\|\tau_y(f*g)-f*g\|_u
=\|(\tau_yf-f)*g\|_u
\leqslant\|\tau_yf-f\|_p\|g\|_q\longrightarrow0\qquad(y\to0).
\end{align*}
若 $p=\infty$，交换 $f,g$ 的角色即可。最后，当 $1<p,q<\infty$ 时，取紧支撑函数序列 $f_j,g_j$，使 $\|f_j-f\|_p\to0$、$\|g_j-g\|_q\to0$。由卷积支撑的包含关系及刚证明的连续性，$f_j*g_j\in C_c$。又
\begin{align*}
\|f_j*g_j-f*g\|_u
\leqslant\|f_j-f\|_p\|g_j\|_q
+\|f\|_p\|g_j-g\|_q\longrightarrow0.
\end{align*}
由\cref{proposition:ch04-cc-dense-czero}，得到 $f*g\in C_0$。
\end{proof}
% SUPPLEMENT: id=v14-ch08-conjugate-convolution-details; parent=ch08-conjugate-convolution; kind=omitted-proof+motivation+hypotheses; origin=independent-derivation
\begin{remark}
\textbf{逐点存在、逼近与端点（重构新增）。}固定 $x$ 后，$y\mapsto f(x-y)$ 的 $L^p$ 范数等于 $\|f\|_p$，所以 Hölder 给 $\int|f(x-y)g(y)|\,\mathrm dy\leqslant\|f\|_p\|g\|_q$，对每个 $x$ 都成立。对差函数同样作逐点估计再取上确界，便得到原证中的一致连续性不等式；无需预先假设卷积连续，也无需先把一致范数换成本质一致范数。共轭条件保证这次 Hölder 可用，而且至少一个指数有限，故可以把平移落在该因子上，使用有限指数的平移连续性。

在两个指数均有限时，可具体取 $f_j=f\chi_{\overline B(0,j)}$、$g_j=g\chi_{\overline B(0,j)}$。可积尾部趋零，给相应的范数收敛；卷积的支撑包含于闭球 $\overline B(0,2j)$，故连续卷积有紧支撑。$\|g_j\|_q\leqslant\|g\|_q$，因此原证乘积误差确实趋零。$C_c$ 的一致极限在无穷远处消失：给定误差，先选一个一致逼近函数，在它的紧支撑外，极限函数的绝对值即被该误差控制。

最后一条结论的严格端点限制不能删掉。例如取非负 $f\in L^1$ 满足 $\int f=1$，取 $g=1\in L^\infty$，则 $f*g=1$，它一致连续有界，却不在 $C_0(\mathbb R^n)$。动机是：共轭可积性提供每一点的配对，而两个有限指数进一步允许同时截去两因子的远处尾部，从而恢复无穷远处的消失性。
\end{remark}
% END SUPPLEMENT: id=v14-ch08-conjugate-convolution-details
% SOURCE p=256; printed=241; object=ch08-general-young-introduction; region=middle
前面的结果已经足够供本书使用；为完整起见，再陈述如下推广。
% SOURCE p=256; printed=241; object=ch08-general-young; number=8.9; region=middle-lower
\begin{proposition}\label{proposition:ch08-general-young}
设 $1\leqslant p,q,r\leqslant\infty$，且 $p^{-1}+q^{-1}=r^{-1}+1$。
\begin{enumerate}
\item\label{item:ch08-general-young-strong}\textbf{Young 不等式的一般形式。}若 $f\in L^p$、$g\in L^q$，则 $f*g\in L^r$，且 $\|f*g\|_r\leqslant\|f\|_p\|g\|_q$。
\item\label{item:ch08-general-young-weak-strong}再设 $p>1$、$q>1$、$r<\infty$。若 $f\in L^p$、$g\in\text{弱 }L^q$，则 $f*g\in L^r$，且 $\|f*g\|_r\leqslant C_{pq}\|f\|_p[g]_q$，其中 $C_{pq}$ 与 $f,g$ 无关。
\item\label{item:ch08-general-young-weak}设 $p=1$、$r=q>1$。若 $f\in L^1$、$g\in\text{弱 }L^q$，则 $f*g\in\text{弱 }L^q$，且
\begin{align*}
[f*g]_q\leqslant C_q\|f\|_1[g]_q,
\end{align*}
其中 $C_q$ 与 $f,g$ 无关。原书此处漏写的因子及无穷端点约定见下注。
\end{enumerate}
\end{proposition}
\begin{proof}
证明 (a) 时固定 $q$。两个端点 $p=1,r=q$ 与 $p=q/(q-1),r=\infty$，分别是\cref{theorem:ch08-young,proposition:ch08-conjugate-convolution}。一般情形由 Riesz--Thorin 插值定理得到；另见\cref{exercise:ch08-direct-young}的直接证明。(b)、(c) 则为\cref{theorem:ch06-weak-kernel}的特例。
\end{proof}
% SUPPLEMENT: id=v14-ch08-general-young-details; parent=ch08-general-young; kind=omitted-proof+source-correction+motivation+hypotheses; origin=independent-derivation
\begin{remark}
\textbf{插值、核条件与原式订正（重构新增）。}先补足 (a)。若 $q=1$，由交换性和基本 Young 不等式直接成立；若 $q=\infty$，指数关系强迫 $p=1,r=\infty$，也是已证端点。设 $1<q<\infty$，记 $q'=q/(q-1)$。固定 $g\in L^q$，线性算子 $Tf=f*g$ 在 $L^1\to L^q$ 与 $L^{q'}\to L^\infty$ 上的范数均不超过 $\|g\|_q$。两个定义在交集上一致，因为都由同一个绝对收敛积分给出。

指数关系给 $1\leqslant p\leqslant q'$。选择 $0\leqslant\theta\leqslant1$ 使
\begin{align*}
\frac1p=(1-\theta)+\frac\theta{q'},\qquad
\frac1r=\frac{1-\theta}{q};
\end{align*}
这两式之差恰给要求的关系。\cref{theorem:ch06-riesz-thorin}于是给算子范数至多 $\|g\|_q$。这里插值所得算子确实仍是原卷积：对 $f\in L^p$，分成 $f_1=f\chi_{\{|f|>1\}}$、$f_2=f\chi_{\{|f|\leqslant1\}}$。由 $|f_1|\leqslant|f|^p$ 及 $|f_2|^{q'}\leqslant|f|^p$，有 $f_1\in L^1$、$f_2\in L^{q'}$。各部分卷积分别几乎处处和处处绝对收敛，因此其和几乎处处等于 $f*g$，与和空间上的插值定义相同。零范数的 $g$ 给零算子，不需对零作幂的运算。

对 (b)、(c) 的有限 $q$ 情形，取 $K(x,y)=g(x-y)$。可取 Borel 代表；每个行、列截面由 $g$ 平移或反射得到，水平集测度不变，故两侧弱 $L^q$ 准范数均为 $[g]_q$。两空间的 Lebesgue 测度均 $\sigma$ 有限，核定理全部假设成立。(b) 的关系给
$1/r=1/p-1/q'>0$，故 $1<p<q'$，并有 $r>p$，正是核定理中已验证的强型范围；$r<\infty$ 排除了临界端点。在 (b) 中 $q=\infty$ 不可能满足指数关系和 $p>1,r<\infty$。在 (c) 中若容许 $q=\infty$，按弱 $L^\infty=L^\infty$、$[g]_\infty=\|g\|_\infty$ 理解，则用基本 Young 不等式，常数取一即可。

原书 (c) 写作 $[f*g]_q\leqslant C_q\|f\|_1$，同时说 $C_q$ 与 $f,g$ 无关，遗漏了 $[g]_q$。这个因子既由核定理给出，也由齐次性所必需。例如在一维取 $f=g=\chi_{[0,1]}$，则 $f*g$ 为非零三角形函数，其弱 $L^q$ 准范数严格正且有限；将 $g$ 换成 $Ag$，左侧乘以任意 $A>0$，而原式右侧不变，矛盾。此处已明确补上该因子。

动机是由两个可解释的端点控制全部中间卷积估计，并允许核只有弱可积性。指数平衡还符合伸缩：若令 $f_t(x)=f(x/t)$、$g_t(x)=g(x/t)$，则 $f_t*g_t(x)=t^n(f*g)(x/t)$。两侧范数分别按 $t^{n+n/r}$ 与 $t^{n/p+n/q}$ 缩放，因此要有对所有尺度统一的非零估计，就必须满足该关系。弱型常数可以依赖指数，但不能吸收任意函数的大小。
\end{remark}
% END SUPPLEMENT: id=v14-ch08-general-young-details
\subfile{section-02-smoothness.tex}
\end{document}
